% sections/08-coevolution.tex — graded coupling edges; headline per protocol v4 claim ledger.
\section{The Co-Evolution Loop: Graded Evidence}
\label{sec:coevolution}

Our second insight concerns coupling between offense and defense. We grade every candidate edge
on a three-level scale --- \emph{documented response} (an explicit, sourced statement or
traceable institutional flow), \emph{plausible coupling} (mechanism-reasonable temporal
adjacency), and \emph{temporal association} (chronology alone) --- and claim only what the
grades support: multiple documented edges exist; field-wide causal co-evolution does not follow
from them.

\subsection{Documented Response Edges}
\textbf{(D1) Competition $\rightarrow$ industrial research program.} Project Zero's Big Sleep
post opens its target rationale by citing the challenge directly: ``Earlier this year at the
DARPA AIxCC event, Team Atlanta discovered a null-pointer dereference in SQLite, which inspired
us to use it for our testing''~\cite{bigsleep2024naptime}. This is the corpus's cleanest causal
edge: named source, stated mechanism, dated.
\textbf{(D2) Maintainer-side variant response.} After Atlantis reported the FTS5 trigram
off-by-one, the SQLite maintainer patched not only the reported function but sibling tokenizer
code paths sharing the pattern, then revised his own patch in a follow-up commit after noting
redundancy~\cite{taasc2024} --- human variant analysis triggered by an AI finding.
\textbf{(D3) Institutional/funding flows.} The loop acquired governance: DARPA and ARPA-H
committed \$200K per team to integrate CRSs into critical software~\cite{tobbuttercup2025};
Team Atlanta recycled half its \$4M prize (\$2M) into continuous autonomous hunting at Georgia
Tech's SSLab with OpenAI contributing API credits~\cite{taafc2025}; OpenSSF chartered a Cyber
Reasoning Systems SIG and hosts OSS-CRS and FuzzingBrain, whose post-competition operation is
reported at 62 vulnerabilities across 26 projects (43 maintainer-confirmed, 36 patched upstream)
for FuzzingBrain and 25 vulnerabilities across 16 projects for OSS-CRS~\cite{openssfaixcc2026}.
These are foundation-relayed, team-reported figures; the funding and governance mechanisms provide evidence that competition-derived capabilities persisted beyond the scored event.
\textbf{(D4) Threat-seeded defensive disclosure.} Google's CVE-2025-6965 disclosure couples its
Threat Intelligence unit to Big Sleep: the flaw was ``known only to threat actors,'' and the
company asserts an imminent exploit was cut off --- ``we believe this is the first time an AI
agent has been used to directly foil efforts to exploit a vulnerability in the
wild''~\cite{bloggoogle2025,nvd6965}. The threat-intelligence premise is unverifiable and the
foiling counterfactual; we cite it as documented defensive process, not verified attack
prevention.

\subsection{Plausible Coupling}
\textbf{(P1)} Attacker-side adoption of agentic operation within roughly twelve months of
defender agentic publicity: Anthropic reports GTG-1002 used Claude Code for 80--90\% of an
espionage campaign's execution --- ``the AI autonomously discovered vulnerabilities in targets
selected by human operators''~\cite{anthropicgtgreport2025,anthropicgtg2025}. Directionality is
unverifiable without adversary telemetry; shared capability tide remains a live alternative.

\subsection{Temporal Association Only}
Benchmark releases trailing industrial demos, press arms-race narratives, and leaderboard
movements populate this tier; per our claim ledger they carry no inferential weight here.

\subsection{Falsification Test and Record Reliability}
The shared-cause objection --- frontier-model releases independently advancing both sides ---
survives globally: it is defeated only for edges D1--D2, where stated mechanisms precede any
model-release timing argument. We therefore do not generalize from the AIxCC--Google hub to the
field. Record reliability itself becomes evidence, once accounting scopes are labeled precisely:
\$29.5M denotes the cumulative program prize commitment announced in 2024, \$8.5M denotes the
final competition pool per DARPA's scoring guide~\cite{darpascoring2025}, and OpenSSF reports
\$30.5M ultimately distributed across both rounds~\cite{openssfaixcc2026} --- distinct
boundaries, not competing totals for one quantity. Separately, DARPA's own closeout corrects
its earlier synthetic-vulnerability count from 70 to 63~\cite{darparesults2025}: an
administrative record correction showing that even authoritative primary sources require
version- and date-aware verification.

\subsection{Summary}
The loop is real where it is documented: a competition generated tooling whose inspiration
seeded an industrial program (D1), whose findings triggered maintainer responses (D2), whose
participants now operate continuously under foundation governance (D3), and whose defender
capability claims now interleave with threat-intelligence-driven disclosure (D4). What the
record cannot yet support is the stronger vernacular claim of an ongoing arms race; that
requires adversary-side measurement this corpus lacks (Sec.~\ref{sec:roadmap}).
